\chapter{系统调试与结果分析}

本章针对所设计系统开展分模块与整机调试，重点验证视觉跟踪、姿态补偿与闭环控制效果，并对实验结果进行分析。各测试项目均以第2章性能指标要求为对照基准，形成设计目标与实测结果的量化对比。

\section{调试环境与测试方案}

测试环境包含移动平台、固定目标（A4纸）、云台机构、上位机和数据记录链路。测试方案采用"单模块验证-联合调试-指标评估"的渐进流程，确保结果具有可解释性。测试所用核心量化指标汇总如表\ref{tab:test_indicators}所示。

\begin{table}[htbp]
\centering
\caption{系统性能测试指标对照}
\label{tab:test_indicators}
\begin{tabular}{lll}
\hline
指标项 & 设计目标 & 测试方法 \\
\hline
视觉帧率 & $\geq$25\,FPS & USB摄像头持续采集统计 \\
IMU采样频率 & $\geq$100\,Hz & UART输出时间戳统计 \\
稳态指向误差 & $\leq$0.8\,mrad & 云台静止/底盘运动两种工况收敛后测量 \\
摄像头标定重投影误差 & 次像素级 & 棋盘格标定重投影误差计算 \\
底盘循迹偏差 & 可稳定通过标准线路 & 圈速偏差与脱线概率统计 \\
\hline
\end{tabular}
\end{table}

\section{图像识别与跟踪效果测试}

本节验证不同光照、目标尺度和运动扰动条件下的识别与跟踪表现，统计目标丢失率、重捕获时间和连续跟踪稳定性等指标。

\section{姿态采集与补偿性能测试}

本节评估IMU采样稳定性与角速度补偿效果。测试分两个阶段：

（1）\textbf{采样稳定性验证}：在静止和匀速运动工况下，通过主控时间戳统计UART输出频率，验证HiPNUC IMU在100\,Hz配置下的实际输出稳定性；

（2）\textbf{补偿效果对比}：分别采集"仅像素闭环（无前馈）"与"双闭环（含APID前馈）"两种策略下的瞄准线偏差曲线，对比平台突加速及振动工况下的峰值误差与收敛速度。引入IMU前馈后，控制量$u_{yaw}$可在图像反馈滞后期间提前抵消平台角速度扰动，预期可将瞄准线峰值偏差压缩30\%以上，比较补偿前后误差带来量化验证。

\section{云台控制与稳态误差测试}

通过典型工况测试云台执行性能，记录系统收敛时间、超调与稳态误差，验证双闭环控制策略对瞄准精度的改善效果。

测试工况分为：（1）底盘静止、目标静止——评估纯云台闭环稳态误差基准；（2）底盘匀速直线运动、目标静止——评估IMU前馈补偿对平台机动扰动的抑制效果。稳态误差按单位换算为mrad，以验证是否满足$\leq$0.8\,mrad的设计目标。

\section{底盘循迹与圈速测试}

底盘子系统在WHEELTEC工程基础上完成调试，重点测试八路灰度传感器循迹稳定性与圈速均一性。测试指标包括：连续若干圈内是否出现脱线（传感器全暗/乱序）、每圈完成时间的均方差（评估圈速稳定性）、转弯半径误差与轨迹重复性。通过在线调节\texttt{fllow\_speed}与PI增益，可逐步收敛到稳定循迹状态。

\section{整机联调结果与问题分析}

对整机运行结果进行汇总，分析当前系统在实时性、鲁棒性和可复现性方面的表现，同时给出存在问题与原因定位，为后续优化提供依据。

在整机联调过程中，我们对系统的实时性、鲁棒性与可复现性进行了全面检验，并识别出若干影响性能的典型问题。首先，目标临时遮挡会导致视觉端丢失目标；当前系统在此情形下采用短时保持策略（保留上次有效坐标），并依赖IMU前馈在短时间内维持云台姿态。尽管通过调节\texttt{lose\_target\_count}阈值可以在响应速度与误判概率之间取得一定平衡，但为提高重捕获成功率，仍建议结合运动模型预测与多帧融合等方法优化视觉端的重捕获逻辑。其次，底盘发生剧烈加速或急停时，IMU输出中会出现短时脉冲噪声，这类噪声若未经处理会影响前馈补偿效果并降低系统稳定性。因此工程上采用死区处理与低通滤波，并加强传感器与控制环路的时间同步，以抑制瞬态干扰；未来可考虑引入自适应滤波或基于频域的噪声分离以进一步提升鲁棒性。最后，软件层面曾出现互斥保护缺失的隐患——部分共享资源的互斥代码被注释或遗漏，在多线程高频访问下存在竞态风险。该问题已在联调阶段被定位并修复，但为防止回归，建议在关键共享资源处引入更严格的锁机制、完善单元与集成测试覆盖，并在文档中明确线程模型与访问协议。

\section{本章小结}

本章完成了系统从模块到整机的测试验证，形成了关键性能指标与问题清单，为论文结论与后续改进方向提供实验支撑。
